\documentclass[12pt,a4paper]{article}
\usepackage[utf8]{inputenc}
\usepackage[french]{babel}
\usepackage[T1]{fontenc}
\usepackage{geometry}
\usepackage{enumitem}
\usepackage{titlesec}
\usepackage{hyperref}
\usepackage{xcolor}
\usepackage{listings}
\usepackage{tcolorbox}
\usepackage{longtable}
\usepackage{amssymb}

\geometry{top=2cm, bottom=2cm, left=2.5cm, right=2.5cm}

% Code listing configuration
\lstset{
    basicstyle=\ttfamily\small,
    breaklines=true,
    frame=single,
    tabsize=4,
    showstringspaces=false,
    captionpos=b
}

\titleformat{\section}
{\normalfont\LARGE\bfseries\color{blue!70!black}}{\thesection}{1em}{}

\titleformat{\subsection}
{\normalfont\Large\bfseries\color{blue!50!black}}{\thesubsection}{1em}{}

\title{\textbf{Guide Complet\\Création d'un Module CRUD\\E-Salle ENSAA}}
\author{Architecture Clean - 5 Couches}
\date{\today}

\begin{document}

\maketitle
\tableofcontents
\newpage

\section{Introduction}

\subsection{Objectif du Document}

Ce guide vous explique \textbf{étape par étape} comment créer un module CRUD complet en suivant l'architecture utilisée dans le projet E-Salle ENSAA.

\textbf{Ce que vous allez apprendre :}
\begin{itemize}
    \item Créer une entité JPA avec Hibernate
    \item Implémenter le pattern Repository
    \item Développer la couche Service avec validation
    \item Créer un Servlet Controller
    \item Construire des vues JSP avec Bootstrap 5
    \item Configurer Hibernate et enregistrer vos entités
    \item Ajouter votre module dans web.xml
    \item Tester votre module complet
\end{itemize}

\textbf{Durée estimée :} 5 jours pour un module complet (voir Section 8.2)

\subsection{Architecture du Projet}

Le projet suit une \textbf{architecture Clean} en 5 couches :

\begin{center}
\begin{tabular}{|c|l|l|}
\hline
\textbf{Couche} & \textbf{Responsabilité} & \textbf{Package} \\
\hline
1. View & Interface utilisateur (JSP) & \texttt{webapp/WEB-INF/views/} \\
\hline
2. Controller & Gestion des requêtes HTTP & \texttt{features/*/controller/} \\
\hline
3. Service & Logique métier & \texttt{features/*/service/} \\
\hline
4. Repository & Accès aux données & \texttt{features/*/repository/} \\
\hline
5. Domain & Entités métier & \texttt{features/*/domain/} \\
\hline
\end{tabular}
\end{center}

\subsection{Principe de Flux}

\begin{center}
\texttt{JSP $\rightarrow$ Servlet $\rightarrow$ Service $\rightarrow$ Repository $\rightarrow$ Database}
\end{center}

\textbf{Règle d'or :} Chaque couche ne communique qu'avec la couche directement en dessous.

\newpage

\section{Étape 1 : Création de l'Entité (Domain)}

\subsection{Vue d'Ensemble}

L'entité représente la \textbf{structure de données} qui sera persistée en base de données. Elle utilise JPA/Hibernate pour le mapping objet-relationnel.

\subsection{Structure du Package}

\begin{tcolorbox}[colback=gray!10]
\texttt{src/main/java/com/esalle/features/salle/domain/Salle.java}
\end{tcolorbox}

\subsection{Annotations JPA Importantes}

\begin{center}
\begin{tabular}{|l|p{10cm}|}
\hline
\textbf{Annotation} & \textbf{Description} \\
\hline
\texttt{@Entity} & Marque la classe comme entité JPA \\
\hline
\texttt{@Table(name="...")} & Spécifie le nom de la table en base \\
\hline
\texttt{@Id} & Clé primaire \\
\hline
\texttt{@GeneratedValue} & Génération automatique de l'ID \\
\hline
\texttt{@Column} & Configuration de la colonne \\
\hline
\texttt{@NotBlank} & Validation : champ non vide \\
\hline
\texttt{@Size} & Validation : taille min/max \\
\hline
\texttt{@Min / @Max} & Validation : valeur numérique \\
\hline
\texttt{@PreUpdate} & Callback avant mise à jour \\
\hline
\end{tabular}
\end{center}

\subsection{Template d'Entité}

\begin{lstlisting}[language=Java, caption=Template Entité Domain]
package com.esalle.features.MODULE.domain;

import javax.persistence.*;
import javax.validation.constraints.*;
import java.time.LocalDateTime;
import java.util.Objects;

@Entity
@Table(name = "nom_table")
public class NomEntite {
    
    @Id
    @GeneratedValue(strategy = GenerationType.IDENTITY)
    private Long id;
    
    @Column(name = "nom_colonne", nullable = false, length = 100)
    @NotBlank(message = "Le champ est obligatoire")
    @Size(max = 100, message = "Ne peut pas depasser 100 caracteres")
    private String champTexte;
    
    @Column(name = "valeur_numerique")
    @NotNull(message = "La valeur est obligatoire")
    @Min(value = 1, message = "Doit etre au moins 1")
    private Integer champNumerique;
    
    @Column(name = "date_creation", nullable = false)
    private LocalDateTime dateCreation;
    
    @Column(name = "date_modification")
    private LocalDateTime dateModification;
    
    // Constructeur par defaut (requis par JPA)
    public NomEntite() {
        this.dateCreation = LocalDateTime.now();
    }
    
    // Constructeur avec parametres principaux
    public NomEntite(String champTexte, Integer champNumerique) {
        this();
        this.champTexte = champTexte;
        this.champNumerique = champNumerique;
    }
    
    // Getters et Setters
    public Long getId() { return id; }
    public void setId(Long id) { this.id = id; }
    
    // ... autres getters/setters ...
    
    // Callback JPA
    @PreUpdate
    public void preUpdate() {
        this.dateModification = LocalDateTime.now();
    }
    
    // Equals et HashCode (bases sur ID uniquement)
    @Override
    public boolean equals(Object o) {
        if (this == o) return true;
        if (o == null || getClass() != o.getClass()) return false;
        NomEntite entity = (NomEntite) o;
        return Objects.equals(id, entity.id);
    }
    
    @Override
    public int hashCode() {
        return Objects.hash(id);
    }
    
    // ToString pour debug
    @Override
    public String toString() {
        return "NomEntite{id=" + id + ", champTexte='" + champTexte + "'}";
    }
}
\end{lstlisting}

\subsection{Exemple Concret : Salle}

\textbf{Champs principaux :}
\begin{itemize}
    \item \texttt{id} : Long (clé primaire)
    \item \texttt{nom} : String (obligatoire, max 100 caractères)
    \item \texttt{numero} : String (obligatoire, unique, max 20 caractères)
    \item \texttt{capacite} : Integer (obligatoire, min 1, max 1000)
    \item \texttt{description} : String (optionnel, max 500 caractères)
    \item \texttt{equipements} : String (optionnel, max 1000 caractères)
    \item \texttt{etage} : Integer (optionnel)
    \item \texttt{batiment} : String (optionnel, max 50 caractères)
    \item \texttt{disponible} : Boolean (défaut true)
    \item \texttt{dateCreation} : LocalDateTime (auto)
    \item \texttt{dateModification} : LocalDateTime (auto)
\end{itemize}

\subsection{Checklist Entité}

\begin{itemize}
    \item[$\square$] Annotation \texttt{@Entity} présente
    \item[$\square$] Nom de table défini avec \texttt{@Table}
    \item[$\square$] Clé primaire avec \texttt{@Id} et \texttt{@GeneratedValue}
    \item[$\square$] Validations Bean Validation sur les champs
    \item[$\square$] Champs \texttt{dateCreation} et \texttt{dateModification}
    \item[$\square$] Constructeur sans paramètres (requis JPA)
    \item[$\square$] Callback \texttt{@PreUpdate} pour dateModification
    \item[$\square$] Méthodes \texttt{equals()}, \texttt{hashCode()}, \texttt{toString()}
    \item[$\square$] Getters/Setters pour tous les champs
\end{itemize}

\newpage

\section{Étape 2 : Création du Repository}

\subsection{Vue d'Ensemble}

Le Repository gère l'accès aux données. Il utilise Hibernate pour les opérations CRUD et les requêtes personnalisées.

\subsection{Structure du Package}

\begin{tcolorbox}[colback=gray!10]
\texttt{src/main/java/com/esalle/features/MODULE/repository/}
\begin{itemize}
    \item \texttt{ModuleRepository.java} (interface)
    \item \texttt{ModuleRepositoryImpl.java} (implémentation)
\end{itemize}
\end{tcolorbox}

\subsection{Interface Repository}

\begin{lstlisting}[language=Java, caption=Template Interface Repository]
package com.esalle.features.MODULE.repository;

import com.esalle.core.base.BaseRepository;
import com.esalle.features.MODULE.domain.NomEntite;

import java.util.List;
import java.util.Optional;

public interface ModuleRepository extends BaseRepository<NomEntite, Long> {
    
    // Methodes de recherche specifiques
    Optional<NomEntite> findByChampUnique(String valeur);
    
    List<NomEntite> findByChampFiltre(String filtre);
    
    List<NomEntite> findAvailable();
    
    List<NomEntite> searchByKeyword(String keyword);
    
    // Methode de verification d'existence
    boolean existsByChampUnique(String valeur, Long excludeId);
}
\end{lstlisting}

\subsection{Implémentation Repository}

\begin{lstlisting}[language=Java, caption=Template Implementation Repository]
package com.esalle.features.MODULE.repository;

import com.esalle.core.base.BaseRepositoryImpl;
import com.esalle.features.MODULE.domain.NomEntite;
import org.hibernate.Session;
import org.hibernate.query.Query;

import java.util.List;
import java.util.Optional;

public class ModuleRepositoryImpl extends BaseRepositoryImpl<NomEntite, Long> 
    implements ModuleRepository {
    
    public ModuleRepositoryImpl() {
        super(NomEntite.class);
    }
    
    @Override
    public Optional<NomEntite> findByChampUnique(String valeur) {
        Session session = sessionFactory.openSession();
        try {
            Query<NomEntite> query = session.createQuery(
                "FROM NomEntite e WHERE e.champUnique = :valeur", 
                NomEntite.class);
            query.setParameter("valeur", valeur);
            NomEntite result = query.uniqueResult();
            return Optional.ofNullable(result);
        } catch (Exception e) {
            throw new RuntimeException("Error finding by unique field", e);
        } finally {
            session.close();
        }
    }
    
    @Override
    public List<NomEntite> findByChampFiltre(String filtre) {
        Session session = sessionFactory.openSession();
        try {
            Query<NomEntite> query = session.createQuery(
                "FROM NomEntite e WHERE e.champFiltre = :filtre ORDER BY e.nom", 
                NomEntite.class);
            query.setParameter("filtre", filtre);
            return query.list();
        } catch (Exception e) {
            throw new RuntimeException("Error finding by filter", e);
        } finally {
            session.close();
        }
    }
    
    @Override
    public List<NomEntite> searchByKeyword(String keyword) {
        Session session = sessionFactory.openSession();
        try {
            Query<NomEntite> query = session.createQuery(
                "FROM NomEntite e WHERE LOWER(e.nom) LIKE LOWER(:keyword) " +
                "ORDER BY e.nom", 
                NomEntite.class);
            query.setParameter("keyword", "%" + keyword + "%");
            return query.list();
        } catch (Exception e) {
            throw new RuntimeException("Error searching", e);
        } finally {
            session.close();
        }
    }
    
    @Override
    public boolean existsByChampUnique(String valeur, Long excludeId) {
        Session session = sessionFactory.openSession();
        try {
            Query<Long> query;
            if (excludeId != null) {
                query = session.createQuery(
                    "SELECT COUNT(*) FROM NomEntite e " +
                    "WHERE e.champUnique = :valeur AND e.id != :excludeId", 
                    Long.class);
                query.setParameter("excludeId", excludeId);
            } else {
                query = session.createQuery(
                    "SELECT COUNT(*) FROM NomEntite e " +
                    "WHERE e.champUnique = :valeur", 
                    Long.class);
            }
            query.setParameter("valeur", valeur);
            return query.uniqueResult() > 0;
        } catch (Exception e) {
            throw new RuntimeException("Error checking existence", e);
        } finally {
            session.close();
        }
    }
}
\end{lstlisting}

\subsection{BaseRepository - Méthodes Héritées}

Le \texttt{BaseRepository} fournit automatiquement :
\begin{itemize}
    \item \texttt{save(T entity)} : Créer ou mettre à jour
    \item \texttt{findById(ID id)} : Trouver par ID
    \item \texttt{findAll()} : Récupérer toutes les entités
    \item \texttt{deleteById(ID id)} : Supprimer par ID
    \item \texttt{existsById(ID id)} : Vérifier l'existence
    \item \texttt{count()} : Compter les entités
\end{itemize}

\subsection{Bonnes Pratiques Repository}

\begin{enumerate}
    \item \textbf{Toujours fermer la session :} Utiliser \texttt{try-finally}
    \item \textbf{Paramètres nommés :} \texttt{:paramName} au lieu de \texttt{?}
    \item \textbf{Typage générique :} Spécifier le type de retour
    \item \textbf{Gestion d'erreurs :} Wrapper les exceptions Hibernate
    \item \textbf{Requêtes ordonnées :} Toujours ajouter \texttt{ORDER BY}
    \item \textbf{Recherche insensible à la casse :} Utiliser \texttt{LOWER()}
\end{enumerate}

\subsection{Checklist Repository}

\begin{itemize}
    \item[$\square$] Interface hérite de \texttt{BaseRepository<T, ID>}
    \item[$\square$] Implémentation hérite de \texttt{BaseRepositoryImpl}
    \item[$\square$] Constructeur appelle \texttt{super(NomEntite.class)}
    \item[$\square$] Méthodes de recherche spécifiques définies
    \item[$\square$] Méthode \texttt{existsByChampUnique} avec excludeId
    \item[$\square$] Sessions Hibernate correctement fermées
    \item[$\square$] Exceptions gérées et wrappées
    \item[$\square$] Requêtes HQL bien formées avec paramètres nommés
\end{itemize}

\newpage

\section{Étape 3 : Création du Service}

\subsection{Vue d'Ensemble}

Le Service contient la \textbf{logique métier} et les validations. Il fait le lien entre le Controller et le Repository.

\subsection{Structure du Package}

\begin{tcolorbox}[colback=gray!10]
\texttt{src/main/java/com/esalle/features/MODULE/service/}
\begin{itemize}
    \item \texttt{ModuleService.java} (interface)
    \item \texttt{ModuleServiceImpl.java} (implémentation)
\end{itemize}
\end{tcolorbox}

\subsection{Interface Service}

\begin{lstlisting}[language=Java, caption=Template Interface Service]
package com.esalle.features.MODULE.service;

import com.esalle.features.MODULE.domain.NomEntite;
import com.esalle.shared.exception.BusinessException;
import com.esalle.shared.exception.DataAccessException;

import java.util.List;
import java.util.Optional;

public interface ModuleService {
    
    // Operations CRUD
    NomEntite create(NomEntite entite) 
        throws BusinessException, DataAccessException;
    
    NomEntite update(NomEntite entite) 
        throws BusinessException, DataAccessException;
    
    Optional<NomEntite> findById(Long id);
    
    List<NomEntite> findAll();
    
    boolean deleteById(Long id);
    
    // Operations de recherche
    Optional<NomEntite> findByChampUnique(String valeur);
    
    List<NomEntite> findByFiltre(String filtre);
    
    List<NomEntite> searchByKeyword(String keyword);
    
    // Operations metier specifiques
    NomEntite toggleStatus(Long id) throws BusinessException;
    
    // Validation
    void validate(NomEntite entite) throws BusinessException;
}
\end{lstlisting}

\subsection{Implémentation Service}

\begin{lstlisting}[language=Java, caption=Template Implementation Service]
package com.esalle.features.MODULE.service;

import com.esalle.features.MODULE.domain.NomEntite;
import com.esalle.features.MODULE.repository.ModuleRepository;
import com.esalle.features.MODULE.repository.ModuleRepositoryImpl;
import com.esalle.shared.exception.BusinessException;
import com.esalle.shared.exception.DataAccessException;
import com.esalle.shared.util.ValidationUtil;
import org.slf4j.Logger;
import org.slf4j.LoggerFactory;

import java.time.LocalDateTime;
import java.util.List;
import java.util.Optional;

public class ModuleServiceImpl implements ModuleService {

    private static final Logger logger = 
        LoggerFactory.getLogger(ModuleServiceImpl.class);
    private final ModuleRepository repository;

    public ModuleServiceImpl() {
        this.repository = new ModuleRepositoryImpl();
    }

    @Override
    public NomEntite create(NomEntite entite) 
        throws BusinessException, DataAccessException {
        try {
            // 1. Validation Bean Validation
            ValidationUtil.validate(entite);
            
            // 2. Validation metier (unicite)
            if (repository.existsByChampUnique(
                entite.getChampUnique(), null)) {
                throw new BusinessException(
                    "Une entite avec cette valeur existe deja");
            }
            
            // 3. Validation metier personnalisee
            validateBusinessRules(entite);
            
            // 4. Dates automatiques
            entite.setDateCreation(LocalDateTime.now());
            entite.setDateModification(LocalDateTime.now());
            
            // 5. Sauvegarde
            return repository.save(entite);
            
        } catch (javax.validation.ValidationException e) {
            throw new BusinessException(
                "Erreur de validation: " + e.getMessage(), e);
        } catch (Exception e) {
            logger.error("Error creating entity", e);
            throw new DataAccessException(
                "Erreur lors de la creation: " + e.getMessage(), e);
        }
    }

    @Override
    public NomEntite update(NomEntite entite) 
        throws BusinessException, DataAccessException {
        try {
            // 1. Validation Bean Validation
            ValidationUtil.validate(entite);

            // 2. Verifier que l'entite existe
            if (entite.getId() == null || 
                !repository.existsById(entite.getId())) {
                throw new BusinessException("Entite non trouvee");
            }
            
            // 3. Validation unicite (exclure l'entite courante)
            if (repository.existsByChampUnique(
                entite.getChampUnique(), entite.getId())) {
                throw new BusinessException(
                    "Une autre entite avec cette valeur existe deja");
            }
            
            // 4. Validation metier
            validateBusinessRules(entite);
            
            // 5. Date de modification
            entite.setDateModification(LocalDateTime.now());
            
            // 6. Mise a jour
            return repository.save(entite);
            
        } catch (javax.validation.ValidationException e) {
            throw new BusinessException(
                "Erreur de validation: " + e.getMessage(), e);
        } catch (Exception e) {
            logger.error("Error updating entity", e);
            throw new DataAccessException(
                "Erreur lors de la mise a jour: " + e.getMessage(), e);
        }
    }
    
    @Override
    public Optional<NomEntite> findById(Long id) {
        return repository.findById(id);
    }
    
    @Override
    public List<NomEntite> findAll() {
        return repository.findAll();
    }
    
    @Override
    public boolean deleteById(Long id) {
        return repository.deleteById(id);
    }
    
    @Override
    public NomEntite toggleStatus(Long id) throws BusinessException {
        Optional<NomEntite> opt = repository.findById(id);
        if (opt.isPresent()) {
            NomEntite entite = opt.get();
            entite.setActif(!entite.getActif());
            return repository.save(entite);
        }
        throw new BusinessException("Entite non trouvee");
    }
    
    @Override
    public void validate(NomEntite entite) throws BusinessException {
        try {
            ValidationUtil.validate(entite);
            validateBusinessRules(entite);
        } catch (javax.validation.ValidationException e) {
            throw new BusinessException(
                "Erreur de validation: " + e.getMessage(), e);
        }
    }
    
    // Methode privee pour regles metier specifiques
    private void validateBusinessRules(NomEntite entite) 
        throws BusinessException {
        // Exemple: verifier qu'une valeur est dans un range
        if (entite.getValeur() < 0) {
            throw new BusinessException(
                "La valeur doit etre positive");
        }
        
        // Exemple: verifier une dependance
        if (entite.getReference() == null) {
            throw new BusinessException(
                "La reference est obligatoire");
        }
    }
}
\end{lstlisting}

\subsection{Checklist Service}

\begin{itemize}
    \item[$\square$] Interface définit toutes les opérations métier
    \item[$\square$] Implémentation injecte le Repository
    \item[$\square$] Validation Bean Validation avec \texttt{ValidationUtil.validate()}
    \item[$\square$] Validation d'unicité dans create et update
    \item[$\square$] Vérification d'existence dans update
    \item[$\square$] Gestion des dates (création/modification)
    \item[$\square$] Exceptions BusinessException pour erreurs métier
    \item[$\square$] Exceptions DataAccessException pour erreurs techniques
    \item[$\square$] Logger pour tracer les erreurs
    \item[$\square$] Méthodes de validation privées pour règles complexes
\end{itemize}

\newpage

\section{Étape 4 : Création du Servlet (Controller)}

\subsection{Vue d'Ensemble}

Le Servlet gère les requêtes HTTP et fait le lien entre la vue (JSP) et le service.

\subsection{Structure du Package}

\begin{tcolorbox}[colback=gray!10]
\texttt{src/main/java/com/esalle/features/MODULE/controller/ModuleServlet.java}
\end{tcolorbox}

\subsection{Template Servlet}

\begin{lstlisting}[language=Java, caption=Template Servlet Controller]
package com.esalle.features.MODULE.controller;

import com.esalle.features.MODULE.domain.NomEntite;
import com.esalle.features.MODULE.service.ModuleService;
import com.esalle.features.MODULE.service.ModuleServiceImpl;
import com.esalle.shared.exception.BusinessException;
import com.fasterxml.jackson.databind.ObjectMapper;

import javax.servlet.ServletException;
import javax.servlet.http.HttpServlet;
import javax.servlet.http.HttpServletRequest;
import javax.servlet.http.HttpServletResponse;
import java.io.IOException;
import java.util.List;
import java.util.Optional;

public class ModuleServlet extends HttpServlet {
    
    private ModuleService service;
    private ObjectMapper objectMapper;
    
    @Override
    public void init() throws ServletException {
        super.init();
        this.service = new ModuleServiceImpl();
        this.objectMapper = new ObjectMapper();
        this.objectMapper.registerModule(
            new com.fasterxml.jackson.datatype.jsr310.JavaTimeModule());
    }
    
    @Override
    protected void doGet(HttpServletRequest request, 
                        HttpServletResponse response) 
            throws ServletException, IOException {
        
        String pathInfo = request.getPathInfo();
        String action = request.getParameter("action");
        
        try {
            if ("add".equals(action)) {
                // Afficher formulaire creation
                handleAddForm(request, response);
            } else if ("edit".equals(action)) {
                // Afficher formulaire modification
                handleEditForm(request, response);
            } else if ("search".equals(action)) {
                // Recherche
                handleSearch(request, response);
            } else if (pathInfo != null && pathInfo.matches("/\\d+")) {
                // Afficher detail par ID
                Long id = Long.parseLong(pathInfo.substring(1));
                handleGetById(id, request, response);
            } else {
                // Liste par defaut
                handleList(request, response);
            }
        } catch (Exception e) {
            request.setAttribute("error", e.getMessage());
            request.getRequestDispatcher(
                "/WEB-INF/views/MODULE/list.jsp")
                .forward(request, response);
        }
    }
    
    @Override
    protected void doPost(HttpServletRequest request, 
                         HttpServletResponse response) 
            throws ServletException, IOException {
        
        String method = request.getParameter("_method");
        
        try {
            // Support DELETE via POST
            if ("DELETE".equals(method)) {
                handleDelete(request, response);
                return;
            }
            
            // Create ou Update
            String action = request.getParameter("action");
            if ("update".equals(action)) {
                handleUpdate(request, response);
            } else {
                handleCreate(request, response);
            }
        } catch (Exception e) {
            request.setAttribute("error", e.getMessage());
            request.getRequestDispatcher(
                "/WEB-INF/views/MODULE/form.jsp")
                .forward(request, response);
        }
    }
    
    // ... Handlers GET, POST, parseEntityFromRequest ...
}
\end{lstlisting}

\subsection{Configuration web.xml}

\begin{lstlisting}[language=XML, caption=Enregistrement du Servlet]
<servlet>
    <servlet-name>ModuleServlet</servlet-name>
    <servlet-class>
        com.esalle.features.MODULE.controller.ModuleServlet
    </servlet-class>
</servlet>
<servlet-mapping>
    <servlet-name>ModuleServlet</servlet-name>
    <url-pattern>/MODULE/*</url-pattern>
</servlet-mapping>
\end{lstlisting}

\subsection{Patterns de Routing}

\begin{center}
\begin{tabular}{|l|l|l|}
\hline
\textbf{URL} & \textbf{Méthode} & \textbf{Action} \\
\hline
\texttt{/MODULE} & GET & Liste toutes les entités \\
\hline
\texttt{/MODULE?action=add} & GET & Formulaire création \\
\hline
\texttt{/MODULE/123?action=edit} & GET & Formulaire modification \\
\hline
\texttt{/MODULE/123} & GET & Afficher détail \\
\hline
\texttt{/MODULE?action=search\&q=...} & GET & Recherche \\
\hline
\texttt{/MODULE} & POST & Créer entité \\
\hline
\texttt{/MODULE/123?action=update} & POST & Modifier entité \\
\hline
\texttt{/MODULE/123?\_method=DELETE} & POST & Supprimer entité \\
\hline
\end{tabular}
\end{center}

\subsection{Checklist Servlet}

\begin{itemize}
    \item[$\square$] Hérite de \texttt{HttpServlet}
    \item[$\square$] Service injecté dans \texttt{init()}
    \item[$\square$] Méthode \texttt{doGet()} pour affichage
    \item[$\square$] Méthode \texttt{doPost()} pour soumission
    \item[$\square$] Support \texttt{\_method=DELETE} pour suppression
    \item[$\square$] Parsing correct du \texttt{pathInfo}
    \item[$\square$] Redirections après POST (pattern Post-Redirect-Get)
    \item[$\square$] Messages de succès/erreur passés via paramètres URL
    \item[$\square$] Gestion des exceptions avec affichage erreur
    \item[$\square$] Méthode \texttt{parseEntityFromRequest()} pour formulaire
\end{itemize}

\newpage

\section{Étape 5 : Création des Vues (JSP)}

\subsection{Vue d'Ensemble}

Les vues JSP affichent l'interface utilisateur. Le projet utilise Bootstrap 5 pour le design.

\subsection{Structure du Package}

\begin{tcolorbox}[colback=gray!10]
\texttt{src/main/webapp/WEB-INF/views/MODULE/}
\begin{itemize}
    \item \texttt{list.jsp} (liste des entités)
    \item \texttt{view.jsp} (détail d'une entité)
    \item \texttt{form.jsp} (formulaire création/modification)
\end{itemize}
\end{tcolorbox}

\subsection{Composants Principaux JSP}

\begin{center}
\begin{tabular}{|l|p{10cm}|}
\hline
\textbf{Fichier} & \textbf{Contenu} \\
\hline
\texttt{list.jsp} & Barre de recherche, filtres, cartes d'entités, boutons actions \\
\hline
\texttt{view.jsp} & Affichage détaillé de tous les champs, actions (modifier, supprimer) \\
\hline
\texttt{form.jsp} & Formulaire avec validation HTML5, boutons annuler/sauvegarder \\
\hline
\end{tabular}
\end{center}

\subsection{Tags JSTL Importants}

\begin{center}
\begin{tabular}{|l|p{8cm}|}
\hline
\textbf{Tag} & \textbf{Utilisation} \\
\hline
\texttt{<c:if>} & Condition simple \\
\hline
\texttt{<c:choose>} & Condition multiple (switch) \\
\hline
\texttt{<c:when>} & Cas dans choose \\
\hline
\texttt{<c:otherwise>} & Cas par défaut \\
\hline
\texttt{<c:forEach>} & Boucle sur collection \\
\hline
\texttt{\$\{variable\}} & Expression EL pour afficher variable \\
\hline
\texttt{\$\{not empty var\}} & Vérifier si variable non vide \\
\hline
\texttt{\$\{param.name\}} & Accéder paramètre URL \\
\hline
\end{tabular}
\end{center}

\subsection{Checklist JSP}

\begin{itemize}
    \item[$\square$] Directive \texttt{<\%@ page contentType="text/html;charset=UTF-8" \%>}
    \item[$\square$] Import des taglibs JSTL core et fmt
    \item[$\square$] Utilisation de Bootstrap 5 pour le design
    \item[$\square$] Icônes Bootstrap Icons
    \item[$\square$] Messages de succès/erreur avec alerts
    \item[$\square$] Formulaires avec validation HTML5
    \item[$\square$] Confirmation JavaScript pour suppression
    \item[$\square$] URLs construites avec \texttt{\$\{pageContext.request.contextPath\}}
    \item[$\square$] Gestion des cas "liste vide"
    \item[$\square$] Responsive design (colonnes Bootstrap)
\end{itemize}

\newpage

\section{Étape 6 : Configuration et Tests}

\subsection{Configuration hibernate.cfg.xml}

Ajouter le mapping de l'entité :

\begin{lstlisting}[language=XML, caption=hibernate.cfg.xml]
<mapping class="com.esalle.features.MODULE.domain.NomEntite"/>
\end{lstlisting}

\subsection{Création du Script SQL}

\begin{lstlisting}[language=SQL, caption=schema.sql]
CREATE TABLE nom_table (
    id BIGINT PRIMARY KEY AUTO_INCREMENT,
    nom VARCHAR(100) NOT NULL,
    numero VARCHAR(20) NOT NULL UNIQUE,
    capacite INT NOT NULL,
    description VARCHAR(500),
    equipements TEXT,
    etage INT,
    batiment VARCHAR(50),
    disponible BOOLEAN DEFAULT TRUE,
    date_creation TIMESTAMP NOT NULL DEFAULT CURRENT_TIMESTAMP,
    date_modification TIMESTAMP,
    INDEX idx_nom (nom),
    INDEX idx_numero (numero),
    INDEX idx_disponible (disponible)
);
\end{lstlisting}

\subsection{Tests Manuels}

\begin{enumerate}
    \item \textbf{Test Liste :} \texttt{http://localhost:8080/app/MODULE}
    \item \textbf{Test Ajout :} \texttt{/MODULE?action=add}
    \item \textbf{Test Détail :} \texttt{/MODULE/1}
    \item \textbf{Test Modification :} \texttt{/MODULE/1?action=edit}
    \item \textbf{Test Recherche :} \texttt{/MODULE?action=search\&q=test}
    \item \textbf{Test Suppression :} Formulaire DELETE sur liste
\end{enumerate}

\subsection{Points de Vérification}

\begin{itemize}
    \item[$\square$] Création fonctionne et redirige vers liste
    \item[$\square$] Validation empêche création avec données invalides
    \item[$\square$] Unicité du champ unique est vérifiée
    \item[$\square$] Modification sauvegarde correctement
    \item[$\square$] Date de modification se met à jour
    \item[$\square$] Suppression fonctionne avec confirmation
    \item[$\square$] Recherche retourne les bons résultats
    \item[$\square$] Messages de succès apparaissent
    \item[$\square$] Messages d'erreur sont clairs
    \item[$\square$] Interface est responsive
\end{itemize}

\newpage

\section{Guide de Développement Rapide}

\subsection{Checklist Complète}

\begin{longtable}{|c|l|l|}
\hline
\textbf{Étape} & \textbf{Fichier} & \textbf{Action} \\
\hline
\endfirsthead
\hline
\textbf{Étape} & \textbf{Fichier} & \textbf{Action} \\
\hline
\endhead
1 & \texttt{domain/NomEntite.java} & Créer entité avec annotations JPA \\
\hline
2 & \texttt{repository/ModuleRepository.java} & Créer interface repository \\
\hline
3 & \texttt{repository/ModuleRepositoryImpl.java} & Implémenter repository \\
\hline
4 & \texttt{service/ModuleService.java} & Créer interface service \\
\hline
5 & \texttt{service/ModuleServiceImpl.java} & Implémenter service \\
\hline
6 & \texttt{controller/ModuleServlet.java} & Créer servlet controller \\
\hline
7 & \texttt{views/MODULE/list.jsp} & Créer vue liste \\
\hline
8 & \texttt{views/MODULE/view.jsp} & Créer vue détail \\
\hline
9 & \texttt{views/MODULE/form.jsp} & Créer vue formulaire \\
\hline
10 & \texttt{web.xml} & Enregistrer servlet \\
\hline
11 & \texttt{hibernate.cfg.xml} & Ajouter mapping entité \\
\hline
12 & \texttt{schema.sql} & Créer table database \\
\hline
13 & Tests & Tester toutes les fonctionnalités \\
\hline
\end{longtable}

\subsection{Ordre de Développement Recommandé}

\begin{enumerate}
    \item \textbf{Jour 1 (2-3h)} : Entité + Repository + SQL
    \item \textbf{Jour 2 (2-3h)} : Service + Validations
    \item \textbf{Jour 3 (3-4h)} : Servlet + Configuration
    \item \textbf{Jour 4 (3-4h)} : Vues JSP (list, view, form)
    \item \textbf{Jour 5 (2-3h)} : Tests et corrections bugs
\end{enumerate}

\subsection{Commandes Git}

\begin{lstlisting}[language=bash, caption=Workflow Git]
# Creer branche feature
git checkout -b feature/MODULE-crud

# Commits incrementaux
git add .
git commit -m "feat(MODULE): Add domain entity"
git commit -m "feat(MODULE): Add repository layer"
git commit -m "feat(MODULE): Add service layer"
git commit -m "feat(MODULE): Add controller and views"

# Push et Pull Request
git push origin feature/MODULE-crud
\end{lstlisting}

\newpage

\section{Guide Étape par Étape Complet}

\subsection{Préparation}

Avant de commencer, assurez-vous d'avoir :
\begin{itemize}
    \item Une branche Git dédiée : \texttt{git checkout -b feature/MODULE-crud}
    \item Le projet compilable sans erreurs
    \item Une base de données PostgreSQL accessible
\end{itemize}

\subsection{Étape 1 : Créer la Structure de Dossiers}

Créez l'arborescence suivante dans \texttt{src/main/java/com/esalle/features/} :

\begin{verbatim}
MODULE/
├── domain/
├── repository/
├── service/
└── controller/
\end{verbatim}

Créez aussi dans \texttt{src/main/webapp/WEB-INF/views/} :

\begin{verbatim}
MODULE/
├── list.jsp
├── view.jsp
└── form.jsp
\end{verbatim}

\subsection{Étape 2 : Créer l'Entité (Domain)}

\textbf{Fichier :} \texttt{domain/NomEntite.java}

Copiez le template de l'Entité (Section 2.4) et adaptez-le à vos besoins.

\textbf{Points clés :}
\begin{itemize}
    \item Définir tous les champs avec leurs types
    \item Ajouter les annotations Bean Validation
    \item Implémenter constructeurs, getters/setters
    \item Ajouter equals(), hashCode(), toString()
\end{itemize}

\subsection{Étape 3 : Enregistrer l'Entité dans Hibernate}

\textbf{Fichier 1 :} \texttt{src/main/resources/hibernate.cfg.xml}

Ajoutez votre entité dans la section Entity mappings :

\begin{lstlisting}[language=XML]
<!-- Entity mappings -->
<mapping class="com.esalle.features.MODULE.domain.NomEntite"/>
\end{lstlisting}

\textbf{Fichier 2 :} \texttt{src/main/java/com/esalle/shared/config/HibernateConfig.java}

Ajoutez dans la méthode \texttt{getConfiguration()} :

\begin{lstlisting}[language=Java]
// Add entity classes here
configuration.addAnnotatedClass(
    com.esalle.features.MODULE.domain.NomEntite.class);
\end{lstlisting}

\subsection{Étape 4 : Créer le Script SQL}

\textbf{Fichier :} \texttt{src/main/resources/db/schema.sql}

Ajoutez votre table :

\begin{lstlisting}[language=SQL]
CREATE TABLE IF NOT EXISTS nom_table (
    id BIGSERIAL PRIMARY KEY,
    champ1 VARCHAR(100) NOT NULL,
    champ2 VARCHAR(20) NOT NULL UNIQUE,
    champ3 INTEGER NOT NULL CHECK (champ3 > 0),
    description VARCHAR(500),
    date_creation TIMESTAMP NOT NULL DEFAULT CURRENT_TIMESTAMP,
    date_modification TIMESTAMP
);

-- Indexes
CREATE INDEX IF NOT EXISTS idx_table_champ1 ON nom_table(champ1);
CREATE INDEX IF NOT EXISTS idx_table_champ2 ON nom_table(champ2);

-- Sample data
INSERT INTO nom_table (champ1, champ2, champ3) VALUES
('Exemple 1', 'EX001', 100),
('Exemple 2', 'EX002', 200);
\end{lstlisting}

\subsection{Étape 5 : Créer le Repository}

\textbf{Fichier 1 :} \texttt{repository/ModuleRepository.java} (interface)

Copiez le template de l'interface Repository (Section 3.3).

\textbf{Fichier 2 :} \texttt{repository/ModuleRepositoryImpl.java} (implémentation)

Copiez le template de l'implémentation Repository (Section 3.4).

\subsection{Étape 6 : Créer le Service}

\textbf{Fichier 1 :} \texttt{service/ModuleService.java} (interface)

Copiez le template de l'interface Service (Section 4.3).

\textbf{Fichier 2 :} \texttt{service/ModuleServiceImpl.java} (implémentation)

Copiez le template de l'implémentation Service (Section 4.4).

\subsection{Étape 7 : Créer le Servlet}

\textbf{Fichier :} \texttt{controller/ModuleServlet.java}

Copiez le template du Servlet (Section 5.3).

\subsection{Étape 8 : Enregistrer le Servlet}

\textbf{Fichier :} \texttt{src/main/webapp/WEB-INF/web.xml}

Ajoutez avant le commentaire \texttt{<!-- Error Servlet -->} :

\begin{lstlisting}[language=XML]
<!-- MODULE Servlet -->
<servlet>
    <servlet-name>ModuleServlet</servlet-name>
    <servlet-class>
        com.esalle.features.MODULE.controller.ModuleServlet
    </servlet-class>
</servlet>
<servlet-mapping>
    <servlet-name>ModuleServlet</servlet-name>
    <url-pattern>/MODULE</url-pattern>
    <url-pattern>/MODULE/*</url-pattern>
</servlet-mapping>
\end{lstlisting}

\subsection{Étape 9 : Créer les Vues JSP}

Créez les 3 fichiers JSP dans \texttt{views/MODULE/} :

\begin{enumerate}
    \item \textbf{list.jsp} : Liste avec recherche et filtres
    \item \textbf{view.jsp} : Affichage détaillé d'une entité
    \item \textbf{form.jsp} : Formulaire création/modification
\end{enumerate}

Utilisez les templates de la Section 6.

\subsection{Étape 10 : Ajouter la Navigation}

\textbf{Fichier :} \texttt{src/main/webapp/WEB-INF/views/common/header.jsp}

Dans la section navigation, ajoutez :

\begin{lstlisting}[language=HTML]
<li class="nav-item">
    <a class="nav-link" 
       href="\$\{pageContext.request.contextPath\}/MODULE">
        <i class="bi bi-icon-name"></i> Nom Module
    </a>
</li>
\end{lstlisting}

\subsection{Étape 11 : Compiler et Tester}

\textbf{Compilation :}
\begin{lstlisting}[language=bash]
mvn clean compile
\end{lstlisting}

\textbf{Tests d'accès :}
\begin{enumerate}
    \item \texttt{http://localhost:8080/app/MODULE} - Liste
    \item \texttt{http://localhost:8080/app/MODULE?action=add} - Formulaire
    \item \texttt{http://localhost:8080/app/MODULE/1} - Détail
\end{enumerate}

\subsection{Étape 12 : Commit Git}

\begin{lstlisting}[language=bash]
git add .
git commit -m "feat(MODULE): Complete CRUD implementation"
git push origin feature/MODULE-crud
\end{lstlisting}

\subsection{Points de Vérification Finaux}

\begin{itemize}
    \item[$\square$] Entité compilable sans erreurs
    \item[$\square$] Entité enregistrée dans Hibernate
    \item[$\square$] Table créée en base de données
    \item[$\square$] Repository implémenté et testé
    \item[$\square$] Service avec validations complet
    \item[$\square$] Servlet enregistré dans web.xml
    \item[$\square$] 3 vues JSP créées
    \item[$\square$] Navigation ajoutée
    \item[$\square$] Tests manuels effectués
    \item[$\square$] Code committé sur Git
\end{itemize}

\newpage

\section{Bonnes Pratiques}

\subsection{Conventions de Nommage}

\begin{center}
\begin{tabular}{|l|l|l|}
\hline
\textbf{Type} & \textbf{Convention} & \textbf{Exemple} \\
\hline
Entité & PascalCase & \texttt{Salle}, \texttt{Filiere} \\
\hline
Table & snake\_case & \texttt{salles}, \texttt{filieres} \\
\hline
Colonne & snake\_case & \texttt{nom\_salle}, \texttt{date\_creation} \\
\hline
Variable Java & camelCase & \texttt{nomSalle}, \texttt{dateCreation} \\
\hline
Méthode & camelCase + verbe & \texttt{findByNumero}, \texttt{createSalle} \\
\hline
Constante & UPPER\_SNAKE\_CASE & \texttt{MAX\_CAPACITY} \\
\hline
Package & lowercase & \texttt{com.esalle.features.salle} \\
\hline
\end{tabular}
\end{center}

\subsection{Gestion des Erreurs}

\begin{itemize}
    \item \textbf{BusinessException :} Erreurs métier (validation, règles)
    \item \textbf{DataAccessException :} Erreurs techniques (DB, Hibernate)
    \item \textbf{ValidationException :} Erreurs Bean Validation
\end{itemize}

\textbf{Principe :} Toujours catcher les exceptions et donner un message clair à l'utilisateur.

\subsection{Validation}

\textbf{3 niveaux de validation :}
\begin{enumerate}
    \item \textbf{Client (JSP) :} Validation HTML5 + JavaScript
    \item \textbf{Bean Validation :} Annotations \texttt{@NotBlank}, \texttt{@Size}, etc.
    \item \textbf{Métier (Service) :} Règles complexes, unicité, dépendances
\end{enumerate}

\subsection{Sécurité}

\begin{itemize}
    \item Utiliser des \textbf{paramètres nommés} dans les requêtes HQL
    \item Encoder les paramètres dans les vues (JSTL le fait automatiquement)
    \item Valider \textbf{côté serveur} (ne jamais faire confiance au client)
    \item Utiliser HTTPS en production
    \item Limiter les tailles de champs (\texttt{maxlength})
\end{itemize}

\subsection{Performance}

\begin{itemize}
    \item \textbf{Fermer les sessions Hibernate :} Toujours dans \texttt{finally}
    \item \textbf{Indexer les colonnes :} Utilisées dans WHERE et JOIN
    \item \textbf{Limiter les requêtes :} Éviter N+1 queries
    \item \textbf{Pagination :} Pour les listes longues
    \item \textbf{Cache :} Utiliser le cache Hibernate pour entités fréquentes
\end{itemize}

\newpage

\section{Troubleshooting}

\subsection{Problèmes Courants}

\begin{longtable}{|p{5cm}|p{9cm}|}
\hline
\textbf{Problème} & \textbf{Solution} \\
\hline
\endfirsthead
\hline
\textbf{Problème} & \textbf{Solution} \\
\hline
\endhead
\texttt{ClassNotFoundException} & Vérifier que l'entité est mappée dans \texttt{hibernate.cfg.xml} ET \texttt{HibernateConfig.java} \\
\hline
\texttt{404 Not Found} & Vérifier le mapping du servlet dans \texttt{web.xml} (url-pattern doit correspondre) \\
\hline
Servlet ne démarre pas & Vérifier que la classe Servlet existe et le package est correct \\
\hline
Table n'existe pas en DB & Exécuter le script \texttt{schema.sql} ou vérifier \texttt{hibernate.hbm2ddl.auto} \\
\hline
\texttt{NullPointerException} dans JSP & Vérifier que l'attribut est bien passé depuis le servlet \\
\hline
Validation ne fonctionne pas & Appeler \texttt{ValidationUtil.validate()} dans le service \\
\hline
Dates non formatées & Ajouter \texttt{JavaTimeModule} à ObjectMapper \\
\hline
Session Hibernate non fermée & Ajouter \texttt{finally \{ session.close(); \}} \\
\hline
Unicité non vérifiée & Implémenter \texttt{existsByChampUnique()} correctement \\
\hline
Redirection en boucle & Vérifier le pattern Post-Redirect-Get \\
\hline
CSS/JS ne charge pas & Vérifier le chemin avec \texttt{\$\{pageContext.request.contextPath\}} \\
\hline
\end{longtable}

\subsection{Debug}

\textbf{Points de vérification :}
\begin{enumerate}
    \item Consulter les \textbf{logs} dans \texttt{logs/localhost\_access\_log.txt}
    \item Utiliser \texttt{System.out.println()} pour tracer (temporaire)
    \item Activer les logs Hibernate : \texttt{hibernate.show\_sql=true}
    \item Vérifier la base de données avec un client SQL
    \item Utiliser les DevTools du navigateur (Network, Console)
\end{enumerate}

\newpage

\section{Exemple Complet : Module Filière}

\subsection{Spécifications}

\textbf{Module :} Filière (Programme d'études)

\textbf{Champs :}
\begin{itemize}
    \item \texttt{nom} : String (obligatoire, unique, max 100 caractères)
    \item \texttt{code} : String (obligatoire, unique, max 10 caractères)
    \item \texttt{cycle} : Enum (PREPARATOIRE, INGENIEUR)
    \item \texttt{annee} : Integer (1 à 5)
    \item \texttt{effectif} : Integer (min 1)
    \item \texttt{coordinateurId} : Long (foreign key vers users)
    \item \texttt{description} : String (optionnel, max 500)
    \item \texttt{actif} : Boolean (défaut true)
\end{itemize}

\subsection{Règles Métier}

\begin{enumerate}
    \item Le \textbf{code} doit être unique
    \item L'\textbf{année} doit être entre 1 et 2 pour PREPARATOIRE
    \item L'\textbf{année} doit être entre 1 et 3 pour INGENIEUR
    \item L'\textbf{effectif} doit être positif
    \item Le \textbf{coordinateur} doit exister et avoir le rôle COORDINATEUR
\end{enumerate}

\subsection{URLs}

\begin{itemize}
    \item \texttt{/filieres} : Liste
    \item \texttt{/filieres?action=add} : Formulaire création
    \item \texttt{/filieres/123} : Détail
    \item \texttt{/filieres/123?action=edit} : Modification
    \item \texttt{/filieres?cycle=PREPARATOIRE} : Filtre par cycle
\end{itemize}

\subsection{Durée Estimée}

\textbf{Total : 5 jours} (selon PROJECT\_MANAGEMENT\_PLAN.tex)
\begin{itemize}
    \item Jour 1 : Entité + Repository
    \item Jour 2 : Service + Validations métier
    \item Jour 3 : Servlet + Configuration
    \item Jour 4 : Vues JSP
    \item Jour 5 : Tests + Intégration
\end{itemize}

\newpage

\section{Conclusion}

\subsection{Récapitulatif}

Ce guide vous a montré comment créer un module CRUD complet en suivant l'architecture Clean du projet E-Salle ENSAA.

\textbf{5 Couches :}
\begin{enumerate}
    \item \textbf{Domain :} Entité JPA avec validations
    \item \textbf{Repository :} Accès aux données avec Hibernate
    \item \textbf{Service :} Logique métier et validations
    \item \textbf{Controller :} Servlet pour gérer les requêtes HTTP
    \item \textbf{View :} JSP avec Bootstrap pour l'UI
\end{enumerate}

\subsection{Prochaines Étapes}

Maintenant que vous maîtrisez la création d'un CRUD, vous pouvez :
\begin{itemize}
    \item Créer les modules \textbf{Filière} et \textbf{Matière} (Développeur 1)
    \item Créer les modules \textbf{Réclamation}, \textbf{Réservation}, \textbf{Emploi du Temps} (Développeur 2)
    \item Ajouter des fonctionnalités avancées (pagination, export, statistiques)
    \item Implémenter les tests unitaires
\end{itemize}

\subsection{Resources}

\begin{itemize}
    \item \textbf{Documentation Hibernate :} \url{https://hibernate.org/orm/documentation/}
    \item \textbf{JSTL Reference :} \url{https://docs.oracle.com/javaee/5/jstl/1.1/docs/tlddocs/}
    \item \textbf{Bootstrap 5 :} \url{https://getbootstrap.com/docs/5.0/}
    \item \textbf{Bean Validation :} \url{https://beanvalidation.org/}
\end{itemize}

\subsection{Contact}

Pour toute question ou clarification, contactez l'équipe projet.

\vspace{1cm}

\begin{center}
\textbf{Bon développement ! 🚀}
\end{center}

\end{document}
